% Separate post-CP36 candidate. No canonical or checkpoint gate change.
\subsection{Irreducible residual recovery at every valuation-one absolute index}

Let $p\ge5$, $q=p^2$, $N=q-1$, and let $K$ be complete
discretely valued of mixed characteristic $(0,p)$ with algebraically
closed residue field and $\mu_p\subset K$. Write
\[
 e=v_K(p)=tp(p-1),\q\quad t\in\mathbf Z_{>0},\quad p\nmid t.
                                                        \tag{B-RD99.1}
\]
Thus $v_p(e)=1$; the integer $t$ need not be prime, bounded,
or prime to $N$. All lower valuations are normalized on their
named top field. A conductor means a last upper break, with no
conductor-plus-one shift. All extensions below are totally ramified.

\begin{lemma}[Actual point fields and their complete terminal bound]
\label{lem:B-RD99-points}
\status{Candidate; complete proof awaiting independent audits.}
Suppose two whole potentially supersingular height-two groups
$M_1,M_2$ have isomorphic actual $p^3$-torsion, irreducible
residual inertia, and different finite special characters.
Let $L/K$ be the common tame good field of degree $N$ and put
$F_j=L(M_i[p^j])$ for $j=2,3$. For any corresponding primitive
points $P^+,P$ with $pP^+=P$, write $T=L(P^+)$ and $E=L(P)$.
Then
\[
 [E:L]=q,\quad[T:E]=q,\quad[T:L]=q^2,\q\quad
 E/L\text{ and }T/E\text{ are Galois}.                 \tag{B-RD99.2}
\]
Writing $W_1,W_2\subset M_2(\mathbf F_p)$ for the actual
first and terminal congruence layers, one has
\[
 W_1\subset W_2,\quad r_1=\dim W_1\ge3,
 \q\quad W_1v=\mathbf F_p^2\quad(v\ne0).               \tag{B-RD99.3}
\]
The exact last upper breaks of $E/L$, $T/E$ and $F_2/L$ are,
respectively, $c_1,c_2,c_1$, where
\[
 c_1=tN(p-1),\q\quad c_2=t(q^2-1)(p-1).
                                                        \tag{B-RD99.4}
\]
Every character of the complete terminal extension $F_3/F_2$
has conductor at most $p^{r_1-2}c_2$.
\end{lemma}
\begin{proof}
We supply the actual-field and coefficient details, rather than
applying the $t=1$ statement outside its scope. A nontrivial
valuation filtration in $M_i[p]$ would give an inertia-stable
line, so irreducibility makes all nonzero first points have
one valuation. The Weierstrass product of their coordinates
has valuation $v_L(p)=eN$; their common valuation is therefore
$e$. The normal wild inertia subgroup fixes a nonzero vector
over $\mathbf F_p$. Its fixed space is inertia-stable and hence
is the whole irreducible module. The remaining irreducible
cyclic tame group embeds in $\mathbf F_q^\times$ and is killed
over $L$. The existence of the common good field itself is the
arbitrary-index Lemma~\ref{lem:B-RD49-inertia}.

The first Newton polygon is the single segment from $(1,eN)$
to $(q,0)$. Its degree-$i$ coefficient for $i<q$ has valuation
at least $(q-i)e$. At a target of valuation $0<Z\le e$ this
is strictly above the inverse chord $(0,Z)$ to $(q,0)$.
The degree-$q$ term is uniquely first at every inverse point.
Consequently a point of order $p^j$, $1\le j\le3$, has
valuation $e/q^{j-1}$ over $L$. In Honda coordinates its
leading multiplication-by-$p$ map is the $q$th power.

We also record the orientation used below. If the special
tangent exponent is $c_i$, first-level leading coefficients
carry tame exponent $c_i+e$. This scalar is not in $\mathbf F_p$,
so their two-dimensional leading space is a scalar multiple
of $\mathbf F_q$. A residual intertwiner is either field-linear,
which identifies $c_1=c_2\pmod N$, or Frobenius-semilinear.
The former would give the same special character. Thus under
our contrary hypothesis the normalized leading translation
parameters satisfy $z_2=z_1^p$. Unique $q$th roots transport
the same relation to the second and third levels. This is the
arbitrary-$e$ orientation proof of B79, now with $e$ as in
(B-RD99.1), not a change of tame generator or of covariant
convention.

All first-level translations are $L$-rational. Therefore
$E/L$ and $T/E$ are Galois translation extensions of degree
at most $q$. Their coordinate denominators give
\[
 v_L(P)=\frac{t(p-1)}p,\q\quad
 v_L(P^+)=\frac{t(p-1)}{p^3},
 \quad [E:L]\ge p,\quad[T:L]\ge p^3.                 \tag{B-RD99.5}
\]
Here $p\nmid t$ is essential. If $[E:L]=p$, then $[T:L]=p^3$
and $T/E$ has all $q$ translations. In $T$ the primitive
coordinate has valuation $A=t(p-1)$, prime to $p$, while
each nonzero translation displaces it by a first-level point,
of valuation $D=ep^3$. Choose one uniformizer of this actual
common field and fixed residue representatives in its base.
The complete uniformizer expansion of B60 gives lower break
$D-A$ and leading displacement $\overline A a_i d_g$ in
model $i$, where $a_i$ is its nonzero initial coefficient and
$d_g$ is the uniformizer displacement. All calculations are
below $v_T(p)=ND$. Thus the two full first-level translation
maps would be scalar proportional. Their relation is Frobenius
on $\mathbf F_q$, which is not scalar, as is seen at $1$ and
an element outside $\mathbf F_p$. This contradiction proves
$[E:L]=q$ for every $P$, and hence $W_1v=\mathbf F_p^2$.

The literal identity $(1+pA+p^2B)^p=1+p^2A\pmod{p^3}$
gives $W_1\subset W_2$ in the same actual coordinates.
Thus terminal matrices give all first-level translations of
every $P^+$ over $E$. This proves $[T:E]=q$ and (B-RD99.2).
In particular $r_1\ge2$.

Suppose $r_1=2$. Evaluation at any nonzero residual vector
is an isomorphism. The stabilizer of a primitive $P^+$ in
the full third-level group then lies in the central terminal
subgroup. It is normal, so this particular $T/L$ is Galois,
$E=F_2$, and $G=\operatorname{Gal}(T/L)$ has order $p^4$.
Its terminal subgroup $U=\operatorname{Gal}(T/E)$ has order
$q$, and taking $p$th powers maps onto $U$, by the same actual
evaluation isomorphism. No abelian or Cartan hypothesis is used.
The two selected coordinates have values
\[
 A=e,\qquad D_1=qA,\qquad D_2=q^2A,\qquad v_T(p)=ND_2.
                                                        \tag{B-RD99.6}
\]
Every nonterminal element displaces by a point of order $p^2$,
every nonidentity terminal element by a nonzero first-level
point. Normalize the actual selected first-level point to
parameter $1$. Leading displacements are
$a_i^q\lambda_i(g)$ and $a_i^{q^2}\lambda_i(h)$, with
$\lambda_2=\lambda_1^p$ and each layer map filling $\mathbf F_q$.

Let $a$ be the minimum lower break of $G$. Every nonidentity
terminal element has larger break: it is a $p$th power, and
the additive first ramification coefficient vanishes on taking
a $p$th power. Fix the full coordinate expansions. Since
$v_p(A)=1$, below $v_T(p)$ the first possible change at break
$b$ is $\min\{A+pb,J_i+b\}$, where $J_i$ is the first
nonzero exponent prime to $p$, allowing $J_i=\infty$.
All later divisible exponents change strictly later than
$A+pb$. Set
\[
 a_0=\frac{D_1-A}{p}=tN(p-1),\q\quad
 j_0=A+(p-1)a_0=D_1-a_0,\qquad j_0\equiv-t\pmod p.
                                                        \tag{B-RD99.7}
\]
If $a<a_0$, cancellation before $D_1$ forces
$J_i=A+(p-1)a$ with a nonzero coefficient in both models.
If $a>a_0$, the change at $D_1$ forces $J_i=D_1-a$.
An impermissible or $p$-divisible claimed ordinary degree
already contradicts its definition. At a terminal break $b>a$
this fixed ordinary term is uniquely earlier than every
divisible term: their difference is $(p-1)(b-a)$ in the
first case, or $(p-1)(b-a)+p(a-a_0)$ in the second.
It therefore gives $J_i+b=D_2$ and makes the two complete
terminal translation maps scalar proportional to $d_h$.
This contradicts their Frobenius orientation.

It remains $a=a_0$. No ordinary degree below $j_0$ can occur.
If $f_i$ denotes its possibly zero coefficient, the change at
$D_1$ for a minimum-break element is
\[
 a_i^q\lambda_i(g)=(A/p)a_i d_g^p+j_0 f_i d_g.
                                                        \tag{B-RD99.8}
\]
Both $A/p$ and $j_0$ are nonzero modulo $p$. If one $f_i$
is nonzero, every terminal break is $b_*=D_2-j_0>a_0$,
with that ordinary term uniquely first. Since
\[
 b_*-\frac{D_2-A}{p}
                  =\frac{(p-1)(D_2-D_1)}p>0,          \tag{B-RD99.9}
\]
the other model at this same break must have its exact same
nonzero first ordinary degree $j_0$. The terminal maps are
again scalar proportional, impossible. If both $f_i=0$,
every nonterminal element has break exactly $a_0$, since a
larger break makes every term change after $D_1$. The leading
ramification map on $G/U$ is then a two-dimensional additive
space, and (B-RD99.8) gives
$\lambda_i(g)=(A/p)a_i^{1-q}d_g^p$ on all of it.
These two full maps are scalar proportional, again impossible.
Thus $r_1\ge3$.

For either extension $E/L$ or $T/E$, normalized coordinate
value is $A=e$, with $v_p(A)=1$, and every nonidentity
translation has the same displacement value $D$, equal to
$qA$ or $q^2A$ respectively. Put $b_0=(D-A)/p$.
If there is one lower break $b<b_0$, cancellation forces an
ordinary degree $J=A+(p-1)b$ and $d_g^{p-1}=\kappa\ne0$
on its $q-1$ nonzero leading coefficients. A polynomial of
degree $p-1$ cannot have those $q-1$ distinct roots.
If $b>b_0$, a unique ordinary term at $D-b$ makes the two
full translation maps scalar proportional. Thus one break
must equal $b_0$.
If there are two lower breaks $a<b$, the first must be
less than $b_0$: otherwise every term changes after $D$ at
the last break. Cancellation at the first break forces
$J=A+(p-1)a$, uniquely first at $b$, and therefore
\[
 J+b=D,\qquad b+(p-1)a=D-A.
                                                        \tag{B-RD99.10}
\]
The last upper break is $(b+(p-1)a)/p=b_0$. This proves
the exact $c_1,c_2$ in (B-RD99.4), including all zero-coefficient
and equal-break boundary possibilities.

Finally choose two actual basis points $P_i$ of the $p^2$
module and primitive basis lifts $P_i^+$. Write $E_i=L(P_i)$,
$T_i=L(P_i^+)$, $D_i=F_2T_i$, and let $U$ be the full terminal
subgroup. The two $E_i$ generate $F_2$, so sums of their cyclic
characters give its exact last conductor $c_1$.
If $H_i,S_i$ are the stabilizers of $P_i,P_i^+$, terminal
evaluation is onto and hence $US_i=H_i$. Consequently
\[
 F_2\cap T_i=E_i,\quad[D_i:F_2]=q,\quad
 [F_2:E_i]=p^{r_1-2},\quad D_1D_2=F_3.               \tag{B-RD99.11}
\]
The group $D_i/E_i$ is elementary abelian: in $H_i$ both
first-layer matrices $A,C$ kill the chosen residual vector,
so $[A,C]$ and the exact $p$th-power matrix also kill it.
Thus commutators and powers lie in $U\cap S_i$, the actual
kernel for this compositum. For upper index $u>c_2$ its
ramification group maps trivially to $T_i/E_i$, and its
intersection with the subgroup over $F_2$ is therefore trivial.
The actual subgroup/quotient Herbrand formula gives the bound
$\psi_{F_2/E_i}(c_2)\le p^{r_1-2}c_2$ for $D_i/F_2$.
The two actual column maps $B\mapsto Be_1,Be_2$ have zero
joint kernel. Their character spaces span the complete terminal
dual, so the same bound holds for every character of $F_3/F_2$.
These are the exact actual-field arguments and ramification
formulas of B97 and \cite[Proposition3.7, Theorem3.10 and
Proposition3.13, complete proofs, pp.~14--16]{ColmezCorpsLocaux};
none depends on $t=1$.
\end{proof}

\begin{lemma}[The scaled cyclotomic alternatives]
\label{lem:B-RD99-cyclotomic}
\status{Candidate; complete proof awaiting independent audits.}
Under (B-RD99.1), the following statements hold for the common
tame extension $L/K$ of degree $N$.
If $\mu_{p^2}\subset L$, then $L(\mu_{p^3})/L$ has degree
$p$ and lower break $tN^2$.
If $\mu_{p^2}\not\subset K$, then $C=K(\mu_{p^3})/K$
is cyclic of degree $p^2$ and its upper breaks satisfy
\[
 \delta_1+p\delta_2\ge tp(p^2-p-1).
                                                        \tag{B-RD99.12}
\]
In particular $\delta_1\le t(p-1)$ implies
$\delta_2\ge t(p^2-p-2+1/p)$.
\end{lemma}
\begin{proof}
Choose a Cohen ring in $\mathcal O_K$ and let $K_0$ be its
fraction field. The uniformizer-basis argument of B98 gives
$[K:K_0]=e$ for this arbitrary perfect residue field. Put
$K_1=K_0(\mu_p)$, $K_2=K_0(\mu_{p^2})$ and
$C_0=K_0(\mu_{p^3})$.
If $\mu_{p^2}\subset L$, its degree over $K$ divides both
$p$ and $N$, so it lies in $K$. Then $K/K_2$ is totally
ramified of degree $t$, prime to $p$, and is tame. The
cyclotomic $C_0/K_2$ has degree $p$ and break $N$, by its
actual uniformizer $\zeta_{p^3}-1$. Tame base change first
by degree $t$ and then degree $N$ gives degree $p$ and break
$tN^2$. The disjointness follows from coprime degrees at
each step; no assumption that $t$ is prime to $N$ is needed.

Now suppose $\mu_{p^2}\not\subset K$. The actual degree
$[K:K_1]=tp$ has $p$-adic order one. The cyclic $p^2$ field
$C_0/K_1$ intersects $K$ in degree at most $p$; its unique
degree-$p$ field is $K_2$, forbidden by hypothesis. Therefore
$K\cap C_0=K_1$, $[C:K]=p^2$ and $[C:C_0]=tp$.
All are actual totally ramified extensions.

For a uniformizer $\pi$ of $K$, let $f$ be its degree-$tp$
Eisenstein polynomial over $K_1$. Its leading derivative term
has valuation
\[
        v_K(tp\pi^{tp-1})=e+tp-1=tp^2-1.
                                                        \tag{B-RD99.13}
\]
For $1\le i<tp$, the other nonzero term $i a_i\pi^{i-1}$
has valuation congruent to $i-1$ modulo $tp$. Indeed
$v_K(a_i)$ is a multiple of $tp$ and
$v_K(i)=tp(p-1)v_p(i)$ is too, including when $p\mid i$.
The residues are distinct and different from the leading
residue $tp-1$. Hence no cancellation removes the least term,
and $d(K/K_1)\le tp^2-1$. Monogenicity and the derivative
formula are exactly the full B98 argument and the original
\cite[Proposition3.18 and complete proof, pp.~17--18]
{ColmezCorpsLocaux}. The degree-$tp$ Eisenstein derivative
for $C/C_0$ has every term of valuation at least $tp-1$,
so $d(C/C_0)\ge tp-1$.

The exact cyclotomic different is still
$d(C_0/K_1)=2p^2(p-1)$. Different transitivity in the two
actual towers, with ramification factors $p^2$ and $tp$,
gives
\[
\begin{split}
 d(C/K)&\ge tp-1+tp\,2p^2(p-1)-p^2(tp^2-1)\\
       &=(p-1)\{tp(p^2-p-1)+p+1\}.
\end{split}                                             \tag{B-RD99.14}
\]
The full trace-dual transitivity and ramification-different
proofs are \cite[Proposition3.17 and Corollary3.19,
pp.~17--18]{ColmezCorpsLocaux}. A cyclic group of order $p^2$
has two positive lower breaks, since its nontrivial $p$th-power
subgroup has strictly later leading ramification coefficient.
Its different is $(p-1)(\delta_1+p\delta_2+p+1)$.
Comparison gives (B-RD99.12) and its consequence. This retains
the rational endpoint and uses neither a finite-residue theorem
nor upper-break integrality.
\end{proof}

\begin{theorem}[Character recovery when the absolute index has valuation one]
\label{thm:B-RD99-irreducible}
\status{Candidate; complete proof awaiting independent audits.}
For every field in (B-RD99.1), actual $p^3$-torsion determines
the finite special characters of whole potentially supersingular
height-two groups with irreducible residual inertia.
\end{theorem}
\begin{proof}
Suppose otherwise and apply the first lemma. On $L$ the actual
determinant is the cyclotomic character. If $W_1$ is traceless,
then $r_1\ge3$ forces $W_1=\mathfrak{sl}_2$ and $r_1=3$.
The determinant is trivial modulo $p^2$, so $\mu_{p^2}\subset L$.
The second lemma gives an actual cyclic degree-$p$ field of
conductor $tN^2$. This exceeds $c_1$, so its intersection with
$F_2$ is $L$. Its actual terminal restriction has conductor
$\psi_{F_2/L}(tN^2)>p^3(tN^2-c_1)$, whereas every terminal
character is bounded by $pc_2$. The difference is
\[
 p^3(tN^2-c_1)-pc_2
                 =tpN(p-1)(p^3-p^2-1)>0,             \tag{B-RD99.15}
\]
a contradiction.

Thus the first-layer trace is nonzero and $\mu_{p^2}$ is
absent from $L$ and $K$. The second lemma supplies a cyclic
degree-$p^2$ field $C/K$. Tame base change by degree $N$
multiplies both its positive upper breaks by $N$, as follows
from the uniformizer-root identity and Herbrand integrals.
Its actual first cyclotomic quotient lies in $F_2$, so
$N\delta_1\le c_1$ and $\delta_1\le t(p-1)$.
Put $d_2=N\delta_2$. Then
\[
 d_2\ge tN(p^2-p-2+1/p)>c_1,\qquad
 q(d_2-c_1)-c_2\ge tN(p^4-3p^3+1)>0.                \tag{B-RD99.16}
\]
The intersection $LC\cap F_2$ is exactly their common
degree-$p$ cyclotomic first layer: $F_2/L$ is elementary
abelian and cannot contain the cyclic $p^2$ extension $LC/L$.
Hence $F_2C/F_2$ is a nonzero actual terminal cyclic quotient.
Above $c_1$ the ramification of the actual compositum lies
over $F_2$ and maps to the final cyclic line until $d_2$.
Its conductor is consequently
\[
 \psi_{F_2/L}(d_2)
   =\psi_{F_2/L}(c_1)+p^{r_1}(d_2-c_1)
   >p^{r_1}(d_2-c_1)>p^{r_1-2}c_2.                 \tag{B-RD99.17}
\]
This contradicts the complete terminal-character bound,
including the determinant direction, proved in the first lemma.
Both possible traces are excluded, proving recovery.
\end{proof}

\begin{remark}[Precise scope and supplementary verification]
\label{rem:B-RD99-scope}
\status{Candidate; no full least-index extension is asserted.}
This result treats only the whole-group irreducible-residual
branch at $v_p(e)=1$. In a perfect-residue application one may
pass to the completed maximal unramified base as in B76; the
absolute index and inertia question are unchanged. The theorem
does not settle reducible residual pairs or mixed whole/filtered
pairs at these indices. In particular it makes no new exact
least-index claim for $(5,40,2)$ or for another first-wild prime.

The complete B79 arbitrary-index Newton/orientation argument,
B97 actual-field and fixed-coefficient proofs, B98 different
proof, and retained Colmez3.7/3.10/3.13/3.17--3.19 source chain
were reread. The new point is the explicit $t$ dependence,
including derivative terms with $p\mid i$, the two zero-coefficient
boundaries and the actual cyclotomic intersections. No new
external-source opening, upper-break integrality theorem or
finite-residue class-field theorem is claimed. The standalone
supplement
\path{code/reduction_descent/verify_valuation_one_irreducible.py}
passed on 20 November 2026 at 06:03:28.889962 UTC in
4.479122429998824 seconds. It checks 15 symbolic identities,
49 sectors at seven primes, including 42 sectors with
$\gcd(t,N)>1$, 103,029 derivative monomials including 1,895
indices divisible by $p$, 98 distinct-residue profiles,
147 critical-boundary checks, 7,893 fixed-ordinary-term
inequalities and 717 two-break Herbrand identities.
The exact rational upper-break census contains 195,496 pairs
and 390,992 terminal comparisons; no integrality of upper
breaks is assumed. The executed program, complete output and
stdout are retained. These symbolic and finite checks supplement
the complete actual-field proof; they infer no formal-group
realization. No mathematical checker failure occurred.
\end{remark}
